% v10. General arithmetic sign theorem for q-pattern (q,q,q,1).
% No basin-positivity assumption is needed for the exclusion step.

\subsection{三つの等しい高分母を持つ core の係数和}
\label{uh10:three-section}

一般の最大分母についての局所符号原理は成立しないが，
三つの高分母が等しい場合には，その三頂点の係数和に符号がある．

\begin{proposition}[三つの等高分母の負の総和]\label{uh10:three-negative}
局所四直線模型の分母が $(q,q,q,1)$，$q\geq3$，であり，
全 row sums と pair products が
\[
 \sum_{j\ne i}p_{ij}=-q_i,
 \qquad p_{ij}p_{ji}=q_iq_j-1
\]
を満たすとする．分母一の頂点を $0$ と書けば
\[
 D_0>0,\qquad \sum_{i\ne0}D_i=-D_0<0.
\]
この結論には basin positivity を仮定する必要はない．
\end{proposition}

\begin{proof}
$m=q-1\geq2$ とする．分母一の row の三 entries は非零の signed divisors of $m$ で，
和は $-1$ である．従って，二つが負で一つが正か，二つが正で一つが負である．

まず $(-a,-b,c)$，$a,b,c>0$，$c=a+b-1$ と書ける場合を考える．
pair products から
\[
 D_0=-1+m\left(\frac1a+\frac1b-\frac1c\right).
\]
$a=1$ なら $c=b$ なので $D_0=m-1>0$ である．
$a>1$ なら $c>b$ なので括弧内は $1/a$ より大きい．
$a\mid m$，$a>0$ より $m/a\geq1$ であり，この場合も $D_0>0$ となる．

残る $(a,b,-c)$，$c=a+b+1$ の場合が存在しないことを示す．
\[
 r=\frac ma,\qquad s=\frac mb,\qquad t=\frac mc,
 \qquad N=q^2-1
\]
と置く．高分母の三頂点における，分母一の頂点へ向かう entries は $r,s,-t$ である．
従って高分母三頂点の内部 entries だけを取り，diagonal を
\[
 A=q+r,\qquad B=q+s,\qquad C=q-t
\]
とすれば，三行の和が全て零である $3\times3$ 行列を得る．
その三つの off-diagonal pair products は全て $N$ である．
二つの向きの三角 cycle products を $X,Y$ と書くと，determinant が零であることから
\[
 X+Y=-\{ABC-N(A+B+C)\},
 \qquad XY=N^3.
\]
従って，実数 $X,Y$ の存在の必要条件は
\begin{equation}\label{uh10:triangle-discriminant}
 \{ABC-N(A+B+C)\}^2\geq4N^3
\end{equation}
である．

一方，$c=a+b+1$ は
\[
 rs-rt-st=\frac{rst}{m}
\]
と同値である．これを用いて展開すると
\[
 ABC-N(A+B+C)=-L+U,
 \qquad L=2q^3-3q,
 \qquad U=r+s-t+\frac{rst}{m}
\]
となる．$c>a,b$ より $t<r,s$ なので $U>0$ であり，
$1\leq r,s,t\leq m$ から
\[
 U\leq2m+m^2=N<L.
\]
従って $(-L+U)^2<L^2$ である．さらに
\[
 4N^3-L^2=3q^2-4>0
\]
なので，\eqref{uh10:triangle-discriminant} に矛盾する．
以上で二つが正の row は排除される．
最後に，四頂点全体の係数和は row-sum 条件から零なので，
高分母三頂点の係数和は $-D_0<0$ である．
\end{proof}
